\section{工程协作与变更流程}
\subsection{项目管理视角}
与 Agent 协同工作的团队需要重新定义“任务状态”：Agent 生成的 patch 需要被纳入待审列表，并与 backlog、issue、任务分解保持一致。Graham 建议用“人类负责需求 + Agent 负责执行”的方式把需求拆成多个小任务，然后把 Agent 的活动写进 sprint board。

\begin{itemize}
    \item 将 Agent 的计划阶段视为“开发任务”，由 PM 或 tech lead 审查。
    \item 在 Agent Act 阶段，保留人工确认机制，确保仅执行经过权限批准的命令。
    \item Verify 阶段把测试报告贴到 issue，以供 QA 回归审查。
\end{itemize}

\subsection{交付与风险沟通}
当 Agent 帮助实现 patch 时，交付状态应明确标识“Agent 草案 + 人工确认”。如果 Agent 失败，则需要自动通知相关 owner，附带日志与 diff，以便及时接手。

 \begin{knowledgebox}{沟通矩阵}
快速记录：Agent 失败 → 通知 owner；Agent 成功但需要 review → 跟进 PR；Agent 执行危险命令 → 请求审批。
\end{knowledgebox}

\begin{warningbox}{不要把 Agent 当成黑盒}
在沟通中要说明 Agent 正在做什么、为什么做，从而减少团队对其信任的恐惧，并确保在失败时能快速接替。
\end{warningbox}

\subsection{本章小结}
\begin{itemize}
    \item 项目管理需要把 Agent 输出作为第一类工作项并纳入看板。
    \item 风险沟通与审批机制是 Agent 大规模部署的前提。
\end{itemize}

